\section{A second-order trace identity}
The following elementary deformation lemma supplies the extra prime
factor that a first-order finite-field sum cannot detect.

\paragraph{Trace lemma.}
Let $p$ be an odd prime, and let $F(z)$ be a monic rational function of
degree $-2$, all of whose numerator and denominator roots are integers.
Assume that its poles have order at most two.  Choose a complete set
$\mathcal R$ of integer representatives modulo $p$, with pairwise
differences prime to $p$, and write every root uniquely as $k+p\ell$,
$k\in\mathcal R$.  At residue $k$, let $N_k,D_k$ be the multisets of these
high roots $\ell$, and put
\[
 h_k=|N_k|-|D_k|,\quad
 U_k(V)=\frac{\prod_{\ell\in N_k}(V-\ell)}
                  {\prod_{\ell\in D_k}(V-\ell)},\quad
 \sigma_k=\sum_{N_k}\ell-\sum_{D_k}\ell.
\]
Suppose:
\begin{enumerate}
\item $h_k\ge-1$ for all $k$;
\item at every residue with $h_k=-1$, the function $U_k$ is the same
      rational function $U$;
\item at every residue with $h_k=0$, one has $U_k(V)=(V-a_k)U(V)$ for an
      integer $a_k$;
\item all nonzero differences of distinct high roots used in the local
      partial fractions are prime to $p$.
\end{enumerate}
Then, as a local rational function of $V$,
\[
       p\sum_{k\in\mathcal R}F(k+pV)\equiv0\pmod {p^2}.       \tag{T2}
\]
In particular, every double-pole and simple-pole coefficient at an
integral high pole of this trace is divisible by $p^2$.

\paragraph{Proof.}
Introduce an independent deformation parameter $P$: replace every root
$k+p\ell$ by $k+P\ell$, retaining its multiplicity and its numerator or
denominator status.  Denote the resulting rational function by $F_P(z)$.
It is monic of degree $-2$ for every $P$.  At $P=0$,
\[
 F_0(z)=\frac{\widetilde Q(z)}{H(z)},\quad
 H(z)=\prod_{k\in\mathcal R}(z-k),\quad
 \widetilde Q(z)=\prod_{k\in\mathcal R}(z-k)^{h_k+1},
 \quad\deg\widetilde Q=p-2.
\]
For $h_k=-1$, write
\[
 F_0(z)=\frac{L_k}{z-k}+L_k A_k+O(z-k),\qquad
 A_k=\sum_{s\ne k}\frac{h_s}{k-s}.
\]
For $h_k=0$, put $L_k=F_0(k)$.  All these quantities are $p$-integral.
Also put $B_k=\sum_{s\ne k}\sigma_s/(k-s)$.
Let $\sigma_*$ be the common value of $\sigma_k$ on the minimal classes;
it is fixed by $U$.  The affine-factor hypothesis implies
$\sigma_k=\sigma_*+a_k$ whenever $h_k=0$.

The ordinary residue theorem applied to the proper rational functions
$F_0$ and
\[
 \left.\frac{\partial F_P}{\partial P}\right|_{P=0}
       =-F_0(z)\sum_s\frac{\sigma_s}{z-s}
\]
gives the exact identities
\[
 \sum_{h_k=-1}L_k=0,\qquad
 \sum_{h_k=-1}L_k(B_k+\sigma_*A_k)
       +\sum_{h_k=0}\sigma_kL_k=0.                    \tag{RI}
\]
The second function has degree at most $-3$ at infinity, so no residue
there is omitted.

Expand the deformed trace locally in $P$.  The minimal classes contribute
\[
 P F_P(k+PV)=L_kU(V)\{1+P(VA_k-B_k)\}+O(P^2),
\]
the classes $h_k=0$ contribute $P L_kU_k(V)+O(P^2)$, and all other classes
contribute $O(P^2)$.  By (RI), the constant term of the trace is zero and
its linear coefficient simplifies exactly to
\[
 U(V)(V+\sigma_*)\,D,\qquad
 D=\sum_{h_k=-1}L_kA_k+\sum_{h_k=0}L_k.                 \tag{D}
\]
Modulo $p$, $H(z)=z^p-z$, so $H'=-1$ and $H''=0$.
At a minimal class, $L_kA_k\equiv-\widetilde Q'(k)$; at a class with
$h_k=0$, $L_k\equiv-\widetilde Q'(k)$; and at every other class this
polynomial derivative vanishes.  Consequently
\[
 D\equiv-\sum_{k\in\mathbb F_p}\widetilde Q'(k)=0\pmod p,
\]
since $\deg\widetilde Q'=p-3<p-1$.
Substituting $P=p$ proves (T2).  The expansions have $p$-integral
coefficients because distinct residue representatives differ by units;
extraction of their local principal parts is valid by the last hypothesis.
This also proves the assertion about pole coefficients.

\section{Application to the Meijer--G rational kernel}
Apply the lemma to
\[
 F_n(z)=\frac{\prod_{u=0}^2(z-(10+u)n)_{10n}}
                  {(z+n)_{16n+1}(z+2n)_{14n+1}},
\]
with $p>29$, $p^2>30n$, and representatives $1,\ldots,p$.
Write $n/p=t+f$.  The hypotheses of the trace lemma hold throughout the
following two fractional intervals, for every $t\ge0$:
\[
              \frac1{10}<f<\frac2{19},\qquad
              \frac{11}{12}<f<\frac{15}{16}.           \tag{S2}
\]
Here is the complete affine root data, before shifting the harmonic
cutoff.  Each interval denotes all integer high roots, with multiplicity.
\begingroup\small\setlength{\arraycolsep}{3.5pt}
\[
\begin{array}{c|ccc|cc}
 f&N_0&N_1&N_2&D_0&D_1\\\hline
 (1/10,2/19)
 &[0,10t]&[t,11t]&[2t,12t]
 &[-17t-2,-t-1]&[-16t-2,-2t-1]\\
 (11/12,15/16)
 &[0,10t+8]&[t+1,11t+9]&[2t+2,12t+10]
 &[-17t-16,-t-2]&[-16t-15,-2t-3]
\end{array}
\]
\endgroup
These are the common high-root intervals in the minimal classes
$h_k=-1$.  For the first row every $h_k=0$ quotient $U_k/U$ is one of
\[
                    V+16t+2,\qquad V+2t+1.
\]
For the second row it is one of
\[
                    V-(2t+1),\qquad V-(10t+9),\qquad V+17t+16.
\]
Which of the last three possibilities occurs depends on the Farey cell;
every possibility is affine, which is all the lemma requires.
At the residue zero one has $h_k=1$, so replacing its representative $0$
by $p$ does not change any minimal or zero class in the table.

The accompanying exact certificate verifies this table on all $5400$
cells and boundary strata of the Farey-$29$ arrangement.  Exactly nine
of its $270$ open abscissa cells satisfy the criterion, and their union is
(S2).  To certify all lifts, the code treats each high-root endpoint as
an affine function of the symbolic integer $t$.  For every zero class it
subtracts its numerator/denominator root multiplicities from the common
minimal data.  The signed interval-boundary multiset reduces to $+1$ at
$a(t)$ and $-1$ at $a(t)+1$, proving $U_k/U=V-a(t)$ identically for every
integer $t\ge0$.  No numerical values of $t$ are used to infer this
identity.  As before, an excluded Farey boundary cannot occur when
$p>29$.  Distinct denominator high poles differ by at most $16n/p+1<p$:
indeed $p^2>30n$ and $p>29$ imply
$16n/p+1<16p/30+1<p$.  Thus the local principal-part
extractions use only unit differences.

\paragraph{The stronger arithmetic readouts.}
Write the partial fractions of $F_n$ as
$\sum_j(a_j/(z+j)^2+b_j/(z+j))$.
For every integral high pole $\ell$, collect the corresponding
coefficients as $A_\ell$ and $B_\ell$.  Its double-pole and simple-pole
coefficients in the trace are $A_\ell/p$ and $B_\ell$.  Therefore (T2)
gives
\[
               A_\ell\in p^3\mathbb Z_{(p)},\qquad
               B_\ell\in p^2\mathbb Z_{(p)}.           \tag{PB}
\]
In particular $Q_n=\sum_\ell A_\ell\in p^3\mathbb Z_{(p)}$.

Choose $\eta=p\lceil2n/p\rceil$.  This lies in $[2n,10n]$: on the first
unlifted interval it is $p$, on the second it is $2p$, and for $t\ge1$ it
is less than $3n$.  The exact readout is unchanged by using this cutoff,
because $F_n',F_n''$ vanish at every integer $2n+1,\ldots,10n$.
If $-j=k+p\ell$, $1\le k\le p$, then
\[
 \left\lfloor\frac{j+\eta}{p}\right\rfloor
                 =\lceil2n/p\rceil-\ell-1,
\]
which is independent of $k$.  Moreover $j+\eta\le27n<p^2$.
The condition $h_k\ge-1$, together with the separation of high poles,
also gives the individual bounds $a_j\in p\mathbb Z_{(p)}$ and
$b_j\in\mathbb Z_{(p)}$.  After multiplying the order-$r$ readout by
$p^{r-1}$, its terms whose harmonic denominators are prime to $p$
therefore vanish modulo $p^2$.
The remaining terms are sums of $A_\ell/p$ and $B_\ell$ with
$p$-integral harmonic weights.  By (PB), they vanish modulo $p^2$ too.
Thus, uniformly throughout both intervals and all their periodic lifts,
\[
 \boxed{\ v_p(Q_n)\ge3,\qquad v_p(P_{3,n})\ge0,\qquad
                              v_p(P_{4,n})\ge-1\ }.    \tag{SG}
\]
This is one additional common prime factor beyond the first-order shell
bound, for both the pure $\zeta(3)$ row and every fixed combination
$A\zeta(3)+B\zeta(4)$.
